\section{Model problem and two-level setting}
\label{sec:model-spaces}
%======================================================================

We introduce the Helmholtz variational problem, the nested coarse and fine
meshes used by the LOD discretization, the coarse neighborhoods needed for
localization, and a projective quasi-interpolation onto the coarse space.
The fine mesh is a computational corrector mesh; no global Helmholtz
reference solution is introduced.

%----------------------------------------------------------------------
\subsection{Helmholtz problem, meshes, and patches}
\label{sec:meshes-patches}
%----------------------------------------------------------------------

Let $\Omega\subset\mathbb R^d$, $d\in\{1,2,3\}$, be a bounded
polyhedral Lipschitz domain with boundary $\Gamma:=\partial\Omega$.
We consider
\begin{equation}
\label{eq:strong-problem}
\left\{
\begin{aligned}
  -\Delta u-\kappa^2u &= f
    && \text{in }\Omega,\\
  u &= 0
    && \text{on }\Gamma_D,\\
  \partial_\nu u &= 0
    && \text{on }\Gamma_N,\\
  \partial_\nu u-\ii\kappa u &= 0
    && \text{on }\Gamma_R.
\end{aligned}
\right.
\end{equation}
The boundary parts are pairwise disjoint up to sets of surface measure
zero and satisfy
\[
  \partial\Omega
  =\overline{\Gamma_D}\cup\overline{\Gamma_N}
   \cup\overline{\Gamma_R},
  \qquad |\Gamma_R|>0.
\]
Either $\Gamma_D$ or $\Gamma_N$ may be empty.  Inhomogeneous data can be
handled by a standard lifting; see, e.g.,
\cite[Chapter~5]{MalqvistPeterseim2020}.

For a measurable $D\subseteq\overline\Omega$, set
\begin{equation}
\label{eq:L2-inner-product}
  (v,w)_D:=\int_D v\,\overline w\,\dd x,
\end{equation}
with the surface-integral interpretation when $D\subseteq\partial\Omega$.
Let
\[
  V:=H^1_{\Gamma_D}(\Omega)
  :=\{v\in H^1(\Omega):v|_{\Gamma_D}=0\}.
\]
For $\kappa\geq\kappa_{\min}>0$, define
\begin{equation}
\label{eq:kappa-norm}
  \norm v_\kappa^2
  :=\norm{\nabla v}_{L^2(\Omega)}^2
    +\kappa^2\norm v_{L^2(\Omega)}^2
\end{equation}
and
\begin{equation}
\label{eq:energy-form}
  b_\kappa(v,w)
  :=(\nabla v,\nabla w)_\Omega+\kappa^2(v,w)_\Omega.
\end{equation}
The corresponding local norm is denoted by $\norm{\cdot}_{\kappa,D}$.
The Helmholtz form is
\begin{equation}
\label{eq:helmholtz-form}
  a(v,w)
  :=(\nabla v,\nabla w)_\Omega
    -\kappa^2(v,w)_\Omega
    -\ii\kappa(v,w)_{\Gamma_R}.
\end{equation}
There is a constant $C_a>0$, independent of $\kappa\geq\kappa_{\min}$,
such that
\begin{equation}
\label{eq:a-continuity}
  |a(v,w)|\leq C_a\norm v_\kappa\norm w_\kappa
  \qquad(v,w\in V).
\end{equation}
For $f\in L^2(\Omega)$, let $F(v):=(f,v)_\Omega$.  We assume that the
continuous variational problem
\begin{equation}
\label{eq:continuous-problem}
  a(u,v)=F(v)\qquad(v\in V)
\end{equation}
is well posed.  Wavenumber-explicit stability results for impedance
problems and their Galerkin discretizations can be found in
\cite{EsterhazyMelenk2012,MelenkSauter2011}.

Let $\mathcal T_H\preceq\mathcal T_h$ be nested conforming simplicial
meshes generated by a fixed boundary-label-preserving refinement rule.
Both families are uniformly shape regular.  The associated conforming
piecewise affine spaces satisfy
\begin{equation}
\label{eq:space-nesting}
  V_H\subset V_h\subset V.
\end{equation}
The coarse space $V_H$ carries the global degrees of freedom.  The fine
space $V_h$ is used only for local correctors, local Riesz problems, and
flux reconstructions.  It may remain fixed over several coarse adaptive
steps so that local matrix factorizations can be cached, but it is refined
whenever the exact-solution certificate detects unresolved fine-scale
information or when nesting would otherwise be lost.

For a coarse entity or a union $S$ of coarse entities, define
\begin{equation}
\label{eq:patch-definition}
  N^0(S):=S,
  \qquad
  N(S):=
  \bigcup_{\substack{T\in\mathcal T_H\\
                     \overline T\cap\overline S\neq\varnothing}}T,
  \qquad
  N^m(S):=N(N^{m-1}(S)).
\end{equation}
Uniform shape regularity gives, for fixed $m$, a bounded-overlap constant
$C_{\mathrm{ol}}(m)$ such that
\begin{equation}
\label{eq:patch-overlap}
  \sum_{T\in\mathcal T_H}
  \norm v_{L^2(N^m(T))}^2
  \leq C_{\mathrm{ol}}(m)\norm v_{L^2(\Omega)}^2.
\end{equation}
The same estimate holds componentwise for gradients.

%----------------------------------------------------------------------
\subsection{Projective quasi-interpolation}
\label{sec:interpolation}
%----------------------------------------------------------------------

The LOD construction requires a local projection onto $V_H$ that
distinguishes coarse information from unresolved components. We use the
following projection-and-averaging operator.

Let $\mathcal N_H$ denote the geometric vertices of $\mathcal T_H$, and
let
\[
  \mathcal N_{H,D}
  :=
  \mathcal N_H\cap\overline{\Gamma_D}
\]
denote the vertices subject to the Dirichlet constraint. For
$z\in\mathcal N_H$, let $\lambda_z$ be the corresponding geometric
nodal hat function. The Dirichlet hat functions are retained in this
geometric partition of unity, so that
\begin{equation}
\label{eq:geometric-partition-unity}
  \sum_{z\in\mathcal N_H}\lambda_z=1.
\end{equation}
This identity will later be used to construct local decompositions of
fine-scale functions.

Let $\Pi_H^{\rm dg}$ denote the elementwise $L^2$-orthogonal projection
onto the discontinuous piecewise affine space $P_1(\mathcal T_H)$.
To convert this discontinuous projection into a conforming function
while preserving the Dirichlet condition, we define the nodal averaging
operator
\[
  E_H:P_1(\mathcal T_H)\to V_H
\]
by
\begin{equation}
\label{eq:nodal-averaging}
  (E_Hq)(z)
  :=
  \begin{cases}
    0,
      & z\in\mathcal N_{H,D},\\[1mm]
    \displaystyle
    \sum_{T\ni z}
    \frac{|T|}
         {\sum_{K\ni z}|K|}
    \,q|_T(z),
      & z\in\mathcal N_H\setminus\mathcal N_{H,D}.
  \end{cases}
\end{equation}
We then define
\begin{equation}
\label{eq:IH}
  I_H:=E_H\circ\Pi_H^{\rm dg}.
\end{equation}
The operator $I_H$ is the local projection-and-averaging construction
described in \cite[Chapter~3]{MalqvistPeterseim2020}. It is projective,
\begin{equation}
\label{eq:IH-projective}
  I_Hv_H=v_H
  \qquad
  \text{for all }v_H\in V_H.
\end{equation}
Projectivity is essential below because it makes $I_H$ a genuine
coarse-scale projection and yields a direct decomposition into coarse
and fine-scale components.

The quasi-local character of $I_H$ is made explicit by the dual
representation
\begin{equation}
\label{eq:IH-dual-representation}
  I_Hv
  =
  \sum_{z\in\mathcal N_H\setminus\mathcal N_{H,D}}
  (v,\rho_z)_\Omega\,\lambda_z,
\end{equation}
where $\rho_z\in P_1(\mathcal T_H)$ satisfies
\[
  \operatorname{supp}\rho_z
  \subseteq
  \overline{N(z)}.
\]
This representation will be used later to control the support enlargement
caused by applying $I_H$ to locally supported functions.

\begin{remark}[Choice of quasi-interpolation]
\label{rem:choice-IH}
The particular choice of $I_H$ is not essential to the abstract LOD
construction. The analysis below requires projectivity, quasi-locality,
local $H^1$-stability, and a local first-order approximation property.
Scott--Zhang-type, coefficient-weighted, and geometry-adapted
quasi-interpolation operators provide alternative choices; see, e.g.,
\cite{ScottZhang1990,EngwerHenningMalqvistPeterseim2019,
PeterseimScheichl2016,HellmanMalqvist2017}.
Different choices may lead to different stability and localization
constants.
\end{remark}

The following estimates quantify the local approximation and stability
properties required in the LOD analysis.

\begin{lemma}[Local approximation and stability of $I_H$]
\label{lem:IH-local-properties}
There exist constants $C_{\rm app}>0$ and $C_I^{\rm st}>0$,
independent of the mesh size and adaptive refinement depth, such that
for every $T\in\mathcal T_H$ and every $v\in V$,
\begin{align}
  H_T^{-1}\norm{v-I_Hv}_{L^2(T)}
  &\leq
  C_{\rm app}
  \norm{\nabla v}_{L^2(N(T))},
  \label{eq:IH-approx}\\
  \norm{\nabla I_Hv}_{L^2(T)}
  &\leq
  C_I^{\rm st}
  \norm{\nabla v}_{L^2(N(T))}.
  \label{eq:IH-stability}
\end{align}
The constants depend only on the spatial dimension, the uniform
shape-regularity bound, and the fixed boundary partition.
\end{lemma}

\begin{proof}
For $z\in\mathcal N_H$, let
\[
  \omega_z
  :=
  \bigcup_{\substack{K\in\mathcal T_H\\ z\in K}}K.
\]
For $z\notin\mathcal N_{H,D}$, the nodal coefficient of $I_Hv$ is
\[
  \alpha_z(v)
  :=
  \frac{1}{|\omega_z|}
  \sum_{K\ni z}|K|\,(\Pi_Kv)(z),
\]
where $\Pi_K$ denotes the $L^2(K)$-projection onto $P_1(K)$.
Weighted Cauchy--Schwarz, scaling on the reference simplex, and the
$L^2$-stability of $\Pi_K$ yield
\[
  |\omega_z|\,|\alpha_z(v)|^2
  \lesssim
  \sum_{K\ni z}|K|\,|(\Pi_Kv)(z)|^2
  \lesssim
  \sum_{K\ni z}\norm{\Pi_Kv}_{L^2(K)}^2
  \leq
  \norm{v}_{L^2(\omega_z)}^2.
\]
Since only the vertices of $T$ contribute to $I_Hv|_T$ and
$\omega_z\subseteq N(T)$ for each such vertex, finite-dimensional
scaling gives
\begin{equation}
\label{eq:IH-local-L2-stability}
  \norm{I_Hv}_{L^2(T)}
  \lesssim
  \norm{v}_{L^2(N(T))}.
\end{equation}

Suppose first that the one-layer vertex patch associated with $T$ does
not meet $\Gamma_D$. Let $c$ be the mean of $v$ over this patch. Since
$I_Hc=c$ on $T$, \eqref{eq:IH-local-L2-stability} and the local
Poincar\'e inequality imply
\[
\begin{aligned}
  \norm{v-I_Hv}_{L^2(T)}
  &\leq
  \norm{v-c}_{L^2(T)}
  +
  \norm{I_H(v-c)}_{L^2(T)}
\\
  &\lesssim
  H_T
  \norm{\nabla v}_{L^2(N(T))}.
\end{aligned}
\]
If the relevant patch meets $\Gamma_D$, the same conclusion follows
with $c=0$ from the corresponding local Friedrichs inequality. This
proves \eqref{eq:IH-approx}.

For the stability estimate, the inverse inequality on $P_1(T)$ and
\eqref{eq:IH-local-L2-stability} give, in the first case,
\[
\begin{aligned}
  \norm{\nabla I_Hv}_{L^2(T)}
  &=
  \norm{\nabla I_H(v-c)}_{L^2(T)}
\\
  &\lesssim
  H_T^{-1}
  \norm{I_H(v-c)}_{L^2(T)}
\\
  &\lesssim
  H_T^{-1}
  \norm{v-c}_{L^2(N(T))}
\\
  &\lesssim
  \norm{\nabla v}_{L^2(N(T))}.
\end{aligned}
\]
The Dirichlet case follows again with $c=0$ and the local Friedrichs
inequality.
\end{proof}

The proof is the local counterpart of the stability and approximation
analysis for projection-and-averaging quasi-interpolation operators;
see \cite[Appendix]{KornhuberPeterseimYserentant2018} and
\cite{ErnGuermond2017}.

%======================================================================
